\centerline{\bf \Large  РЕШЕНИЯ и КРИТЕРИИ}

\centerline{\bf \Large  6 КЛАСС}

\problem{\textbf{1}} Выпишем все возможные варианты трёх натуральных чисел, произведение которых равно 12.

$12\cdot1\cdot1=12$\\
$6\cdot2\cdot1=12$\\
$4\cdot3\cdot1=12$\\
$3\cdot2\cdot2=12$\\

Очевидно, что сумма натуральных чисел - это число натуральное, а нам нужно получить 0. Поэтому используем два отрицательных целых числа, чтобы произведение оставалось положительным. Единственный вариант:

4, -3, -1.

Ответ: 4, -3, -1.
\par

\textbf{Критерии.} \\
\textit{7 баллов} $-$ Полностью верное решение\\
\textit{5 баллов} $-$ Правильный ответ без объяснений\\
\textit{0-3 балла} $-$  За все остальные попытки решить задачу\\



\problem{\textbf{2}} Пусть скорость Арслана $X$, скорость Батыра $Y$, а длина одного круга $S$.

Составим два уравнения:

$4X-4Y=9S$

$7X+7Y=18S$

Решив которое получим следующее равенство:

$X=15Y$

Ответ: в 15 раз быстрее.
\par

\textbf{Критерии.} \\
\textit{0-4 баллов} $-$ Решение пункта А\\
\textit{0-1 балла} $-$ Составление условия пункта Б \\
\textit{0-2 баллов} $-$ Решение пункта Б\\
\textit{0 баллов} $-$ Отсутствие решения пункта А (сразу НОЛЬ! баллов)\\
\textit{0 баллов} $-$ Любой ответ без объяснений\\

\problem{\textbf{3}} Так как прямоугольный лист может быть установлен как вертикально, так и горизонтально, то у нас возможны две ситуации. Пусть $X$ - это количество вертикальных разрезаний, а $Y$ - это количество горизонтальных. Тогда имеем $295	\leqslant 20X 	\leqslant 325$ и $385	\leqslant 25Y 	\leqslant 405$. 

$1!) X=15, Y=16$

$2!) X=16, Y=16$

Или $295	\leqslant 25X 	\leqslant 325$ и $385	\leqslant 20Y 	\leqslant 405$.

$3!) X=12, Y=20$

$4!) X=13, Y=20$

Количество частей на один больше количества разрезаний. А количество прямоугольников - это произведение количества частей после вертикальных разрезаний и количества частей после горизонтальных разрезаний.

$1) 16\cdot17=272$

$2) 17\cdot17=289$

$3) 13\cdot21=273$

$4) 14\cdot21=294$

Ответ: $272, 273, 289, 294$

\textbf{Критерии.} \\
\textit{7 баллов} $-$ Полностью верное решение\\
\textit{По 1 баллу} $-$ За каждое нахождение 1!,2!,3! и 4! \\ 
\textit{1 балл} $-$ Количество частей на 1 больше количества разрезаний\\
\textit{1 балл} $-$ Количество прямоугольников - это произведение количества вертикальных частей на горизонтальных\\
\textit{1 балл} $-$ Указано что возможны две ситуации в зависимости от изначального расположения листа\\
\textit{1 балл} $-$  За четыре правильных ответа без объяснений\\ 
\textit{0-1 балл} $-$ За все остальные попытки решить задачу\\
\textit{0 баллов} $-$ Меньше четырех правильных ответов без объяснений\\
\par

\centerline{\bf }
\centerline{\bf }
\centerline{\bf }
\centerline{\bf }
\centerline{\bf }
\centerline{\bf }
\centerline{\bf }
\centerline{\bf }

\problem{\textbf{4}} Оценка: За 14 часов максимум покрасит в синий $14\cdot7=98$ овец, а нужно покрасить 100. Заметим, что каких бы овец не перекрашивал Очир, количество синих изменяется следующим образом:

+7, +5, +3, +1, -1, -3, -5, -7. 

То есть изменяется на нечётное количество. Значит, за 15 часов изменится также на нечётное количество (как сумма нечётного количества нечётных слагаемых). А изменение должно быть с 0 до 100, то есть чётным. Значит, и за 15 часов тоже не сможет перекрасить в синий все 100 овец. 

Пример: Первые 14-ть часов перекрашиваем 98 овец с красного в синий. За 15-ый час перекрашиваем 6 синих овец в красный, а одну красную в синий. И за 16-ый час перекрашиваем 7 красных в 7 синих. Таким образом за 16 часов все овцы станут с синей меткой.\par

\textbf{Критерии.} \\
\textit{3 балла} $-$ Правильный пример на 16 часов\\
\textit{1 балл} $-$ Доказательство невозможности за 14 часов\\
\textit{0-3 баллов} $-$ Доказательство невозможности за 15 часов, где 3 балла это правильное доказательство через чётность\\
\textit{0-1 балл} $-$  За все остальные попытки решить задачу\\ 
\textit{0 баллов} $-$ Любой ответ без объяснений\\


\problem{\textbf{5}} Найдём количество чисел, которые делятся на 5: 200/5=40

Количество чисел, которые делятся на 25: 200/25=8

Количество чисел, которые делятся на 5, но не делятся на 25: 40-8=32

Которые можем разбить на 16 пар и покрасить каждую в различные цвета.

Каждое число, которое делится на 25, покрасим также в различные цвета.

Максимум можем покрасить в 16+8=24 цвета все числа, чтобы условие задания выполнялось. Следовательно, Очир не может этого сделать.\par

\textbf{Критерии.} \\
\textit{7 баллов} $-$ Полностью верное решение\\
\textit{1 балл} $-$ Найдено количество чисел, которые делятся на 5\\
\textit{1 балл} $-$ Найдено количество чисел, которые делятся на 25\\
\textit{1 балл} $-$ Найдено количество чисел, которые делятся на 5, но не делятся на 25\\ 
\textit{1 балл} $-$ Разбиение 32-ух чисел на 16 пар и раскрашивание их в 16 различных цветов\\
\textit{1 балл} $-$ Раскрашивание чисел, которые делятся на 25 в 8 различных цветов\\
\textit{0-1 балл} $-$ За все остальные попытки решить задачу\\
\textit{0 баллов} $-$ Любой ответ без объяснений\\